\documentclass[11pt,a4paper]{article}
\usepackage{fontspec,amsmath,amssymb,amsthm,microtype}
\setmainfont{lmroman10-regular.otf}[BoldFont=lmroman10-bold.otf,ItalicFont=lmroman10-italic.otf,BoldItalicFont=lmroman10-bolditalic.otf]
\setmonofont{lmmono10-regular.otf}[Scale=MatchLowercase]
\usepackage[margin=25mm]{geometry}
\usepackage[unicode,hidelinks]{hyperref}
\setlength{\emergencystretch}{3em}
\setlength{\parskip}{0.4em}
\setlength{\parindent}{0pt}
\allowdisplaybreaks[2]
\raggedbottom
\title{\texorpdfstring{Gibbs finiteness and the exact phase moment from strong convexity}{Gibbs finiteness and the exact phase moment from strong convexity}}
\author{LangevinLean mathematical exposition}
\date{}
\begin{document}\maketitle
\noindent Source module: \texttt{Kinetic\allowbreak{}Langevin/\allowbreak{}Hamiltonian\allowbreak{}Moments.\allowbreak{}lean}. The accompanying HTML pairs this mathematical exposition with the exact Lean statements and records the scope of the formalization.


The weak Hamiltonian limit needs finite mass and an integrable velocity-plus-force quadratic envelope. This module derives both for the actual exponential Hamiltonian density in finite dimensions. It factors that density into a radial Gaussian velocity weight and the potential Gibbs weight, uses existing direct Gaussian-domination estimates, and transfers position moments to force moments by Lipschitz growth.


\section{\texorpdfstring{Factor the actual Hamiltonian density}{Factor the actual Hamiltonian density}}


Because \(H(v,x)=f(x)+\|v\|^2/2\), the positive exponential weight factors as \(e^{-H(v,x)}=e^{-\|v\|^2/2}e^{-f(x)}\). Therefore its product-Haar weighted measure is the product of the separately weighted velocity and position measures. Lean applies the theorem for products of measures with densities, combines the two nonnegative real weights through \texttt{ENNReal.\allowbreak{}of\allowbreak{}Real\_\allowbreak{}mul}, and checks the exponent identity. Normalization has not been assumed or imposed at this step.



{\raggedright\noindent Lean declarations: \texttt{hamiltonian\allowbreak{}Gibbs\_\allowbreak{}factor}.\par}

\section{\texorpdfstring{Derive finite mass of both factors and their product}{Derive finite mass of both factors and their product}}


The velocity weight is the radial Gaussian with precision \(1/2\), whose ordinary integral is finite. The position weight has finite mass by the previously proved strong-convexity Gaussian domination: a supporting quadratic lower bound for \(f\) makes \(e^{-f}\) integrable. Lean converts these finite ordinary integrals to finite extended density integrals, instantiates finite weighted measures, and uses the exact factorization to prove finiteness of the full phase measure.



{\raggedright\noindent Lean declarations: \texttt{hamiltonian\allowbreak{}Velocity\_\allowbreak{}finite}, \texttt{hamiltonian\allowbreak{}Position\_\allowbreak{}finite}, \texttt{hamiltonian\allowbreak{}Gibbs\_\allowbreak{}finite}.\par}

\section{\texorpdfstring{Derive velocity and position second moments before normalization}{Derive velocity and position second moments before normalization}}


A quadratic factor is absorbed by radial Gaussian decay, so \(\int\|v\|^2e^{-\|v\|^2/2}\,dv<\infty\). The same strong-convexity tail estimate gives \(\int\|x\|^2e^{-f(x)}\,dx<\infty\). Lean uses the ordinary-integral characterization of integration under a density, checks that each exponential density is finite and measurable, and invokes the existing Gaussian-tail estimates. These are actual moments of the unnormalized factors, not moment contracts attached to a stationary process.



{\raggedright\noindent Lean declarations: \texttt{hamiltonian\allowbreak{}Velocity\_\allowbreak{}norm\_\allowbreak{}sq\_\allowbreak{}integrable}, \texttt{hamiltonian\allowbreak{}Position\_\allowbreak{}norm\_\allowbreak{}sq\_\allowbreak{}integrable}.\par}

\section{\texorpdfstring{Transfer position moments to force moments}{Transfer position moments to force moments}}


Gradient Lipschitz continuity implies \(\|\nabla f(x)\|\le M\|x\|+\|\nabla f(0)\|\). Squaring and using \((a+b)^2\le2a^2+2b^2\) gives \(\|\nabla f(x)\|^2\le2M^2\|x\|^2+2\|\nabla f(0)\|^2\). The right-hand side is integrable under the finite position Gibbs weight by its second moment. Lean derives the force growth inequality from the Lipschitz difference bound and triangle inequality, then applies domination by this explicit integrable envelope.



{\raggedright\noindent Lean declarations: \texttt{hamiltonian\allowbreak{}Position\_\allowbreak{}force\_\allowbreak{}sq\_\allowbreak{}integrable}.\par}

\section{\texorpdfstring{Assemble the precise phase envelope needed by the weak limit}{Assemble the precise phase envelope needed by the weak limit}}


On the factored finite phase measure, \(\|v\|^2\) is integrable by pulling the velocity moment through the first coordinate, and \(\|\nabla f(x)\|^2\) is integrable by pulling the force moment through the second coordinate. Their sum is exactly the envelope appearing in the rejection theorem. Lean rewrites the phase measure by its proved factorization, supplies finiteness of both factors, and adds these two coordinate integrability results. Thus every tail premise in the later Hamiltonian stationarity limit follows from positive strong convexity and a Lipschitz gradient.



{\raggedright\noindent Lean declarations: \texttt{hamiltonian\allowbreak{}Gibbs\_\allowbreak{}phase\_\allowbreak{}moment}.\par}

\end{document}
